\documentclass[12pt]{article}
\usepackage[utf8]{inputenc}
\usepackage[T1]{fontenc}
\usepackage{hyperref}
\title{Sentinel Language Resources for K-12 Curriculum Development}
\author{Compiled by Claude Sonnet 5.5 (claude-sonnet-5-5) for Andr\'as Kornai}
\date{}
\begin{document}
\maketitle

\section{Introduction}
Sentinelese (ISO 639-3: std; Glottocode sent1241) is spoken by the people of North Sentinel Island in the Andaman Islands, India. The community has deliberately refused contact with outsiders. Searches found no word lists, recordings, grammar or any other linguistic documentation. Its classification is unresolved. An Andamanese or Ongan connection has been suggested, but only on indirect evidence, and Onge speakers brought to the island reportedly could not understand the islanders. A K-12 curriculum in this language cannot be built from any existing resource. Studying or recording the language without the community's consent would also raise serious ethical and legal problems.

\section{Primary Educational Resources}
None were found. There is no primer, textbook or school programme, and no native-language material exists in any accessible form.

\section{Linguistic and Reference Materials}
\begin{itemize}
\item Wikipedia, ``Sentinelese language'' (a general overview that says the language is undocumented): \url{https://en.wikipedia.org/wiki/Sentinelese_language}
\item Rosetta Project, entry for std (reports that no data is available): \url{https://rosettapanglossia.longnow.org/index.php?title=Std}
\item The Language Closet, ``What do we know about the Sentinelese language?'' (a popular blog summary): \url{https://thelanguagecloset.com/2021/10/23/what-do-we-know-about-the-sentinelese-language/}
\end{itemize}
Searches also returned several content-farm pages. Some of them make specific claims about the language, such as sound inventories and vocabulary sizes. These claims are not backed by fieldwork, so none of those pages are listed here.

\section{Religious and Community Resources}
None were found. The community does not engage with outsiders, and no community organisation represents it.

\section{Digital and Audio Resources}
No audio, video or corpus exists in the search results.

\section{Curriculum Development Context}
The reference description says the language is not known to be taught in schools. No documentation exists on which to base a curriculum. Any future work would need to respect the community's refusal of contact and the protections that Indian law gives it. Related Andamanese languages are better documented, but they must not be mixed up with Sentinelese.
\end{document}